\documentclass[../main]{subfiles}
\begin{document}

\section{Vistas}

\subsection{Vistas ISO(E)}
Como ya sabemos las 3 vistas de un objeto son el alzado, la planta y el perfil. Pero resulta que hay 2 sistemas diferentes de representar estas vistas, el europeo y el americano.
\begin{figure}[h]
    \centering
\includegraphics[]{images/07/vistas-americano.jpg}
    \caption{Sistema Americano}
    \label{fig:my_europeo}
\end{figure}


Veamos las diferencias. Lo primero, se denomina las vistas principales de un objeto, a las proyecciones ortogonales (perpendiculares) del mismo sobre 6 planos, dispuestos en forma de cubo. También se podría definir las vistas como, las proyecciones ortogonales de un objeto, según las distintas direcciones desde donde se mire. Si situamos un observador según las seis direcciones indicadas por las flechas, obtendríamos las seis vistas posibles de un objeto.
\begin{figure}[h]
    \centering
\includegraphics[]{images/07/vistas-europeo.jpg}
    \caption{Sistema Americano}
    \label{fig:my_sistemaamericano}
\end{figure}
El desarrollo del cubo de proyección, nos proporciona sobre un único plano de dibujo, las seis vistas principales de un objeto, en sus posiciones relativas.         

 Con el objeto de identificar, en que sistema se ha representado el objeto, se debe añadir el símbolo que se puede apreciar en las figuras, y que representa el alzado y vista lateral izquierda, de un cono truncado, en cada uno de los sistemas.

\includegraphics[width=0.9\textwidth]{images/07/sistema-americano-europeo.jpg}


\newpage
\subsection{Copiar la siguiente figura.}
\includegraphics[]{images/07/dibujo1}

\end{document}